% Gantry model augmentation: setup for review by Jan Hoekstra.
% Notation follows Hoekstra et al. (2025), "Learning-based model augmentation with LFRs".
% Compile: pdflatex jan-augmentation-writeup.tex . No code listings on purpose.
\documentclass[11pt,a4paper]{article}

\usepackage{amsmath, amssymb, bm}
\usepackage[margin=2.2cm]{geometry}
\usepackage{xcolor}
\usepackage{parskip}
\usepackage{enumitem}
\setlist{nosep, leftmargin=1.4em}

\newcommand{\hl}[1]{\colorbox{yellow!35}{$\displaystyle #1$}}
\newcommand{\xb}{\tilde{x}}   % baseline state
\newcommand{\xa}{\bar{x}}     % added state
\newcommand{\xh}{\hat{x}}     % full model state
\newcommand{\yh}{\hat{y}}
\newcommand{\Cbar}{\bar{C}_d}
\newcommand{\fb}{f_{\mathrm{base}}}
\newcommand{\hb}{h_{\mathrm{base}}}
\newcommand{\fa}{f_{\mathrm{aug}}}
\newcommand{\ga}{g_{\mathrm{aug}}}
\newcommand{\ha}{h_{\mathrm{aug}}}

\begin{document}

\begin{center}
  {\large\bfseries Gantry model augmentation: setup for review}\\[2pt]
  Dirk van den Berg \quad\textbullet\quad \today \\[2pt]
  {\small Notation follows Hoekstra et al.\ (2025). Equation numbers in parentheses refer to that paper.}
\end{center}

% ---------------------------------------------------------------------------
\section*{Notation}
\begin{tabular}{@{}ll@{}}
$\xb \in \mathbb{R}^6$ & baseline state $[q_1,q_2,q_3,\dot q_1,\dot q_2,\dot q_3]$, logical coords (Eq.~2)\\
$\xa \in \mathbb{R}^2$ & added states, targeting $[\delta_a,\dot\delta_a]$ (Sec.~2)\\
$\xh = [\xb^\top\ \xa^\top]^\top \in \mathbb{R}^8$ & full model state (Eq.~3)\\
$u,\ y \in \mathbb{R}^3$ & stage input (forces $F_1,F_2,F_y$), stage output (positions $X_1,X_2,Y$) \\
$Y = q_3$ & scheduling variable \\
$\fb,\hb$ & baseline state-transition / output; \ $\fa,\ga,\ha$ augmentation functions (Table~1)\\
\end{tabular}

% ---------------------------------------------------------------------------
\section{System: the data-generating truth \quad (Eq.~1)}

An 8-state gantry with a hidden tuned-mass absorber. Coordinates
$q = [X,\ \Theta,\ Y,\ \delta_a]^\top$; the hidden mass $m_a$ sits at
$Y{+}L_0{+}\delta_a$ on the cross-beam with spring $k_a$ and damper $c_a$. The
continuous-time equations of motion are
\begin{equation}
  M(Y,\delta_a)\,\ddot q = \begin{bmatrix} P^\top u \\ 0 \end{bmatrix}
    - C_4\,\dot q - K_4\,q ,
\end{equation}
\begin{equation}
  \small
  M(Y,\delta_a) =
  \begin{bmatrix}
    m_{\mathrm{t}} & s & 0 & 0 \\
    s & J & -(m_h{+}m_a)d & -m_a d \\
    0 & -(m_h{+}m_a)d & m_h{+}m_a & m_a \\
    0 & -m_a d & m_a & m_a
  \end{bmatrix},
  \quad
  C_4 = \begin{bmatrix} C & 0 \\ 0 & c_a \end{bmatrix},
  \quad
  K_4 = \begin{bmatrix} K & 0 \\ 0 & k_a \end{bmatrix},
\end{equation}
with (here $m_h = m_{h,\mathrm{rigid}}$, so $m_h{+}m_a = 10.1$):
\begin{align*}
  m_{\mathrm{t}} &= m_1{+}m_2{+}m_b{+}m_h{+}m_a, &
  s &= (m_1{-}m_2)\tfrac{L_b}{2} - (m_h{+}m_a)Y - \hl{m_a(L_0{+}\delta_a)}, \\
  J &= J_b{+}J_h{+}(m_1{+}m_2)\tfrac{L_b^2}{4} + (m_h{+}m_a)d^2 + m_h Y^2 + \hl{m_a(Y{+}L_0{+}\delta_a)^2}. &&
\end{align*}
Discretized (RK4) this is the truth $x_{k+1} = f(x_k,u_k)$, $y_k = h(x_k)$ of Eq.~1.
The highlighted terms are the absorber's coupling: through the mass matrix it enters the
$X$, $\Theta$, $Y$ accelerations (it is not a single-axis spring). Parameters:
$m_a = 0.10\,m_h = 1.01$ kg, $f_a = 150$ Hz with $k_a = m_a(2\pi f_a)^2$,
$\zeta_a = 0.05$ with $c_a = 2\zeta_a\sqrt{k_a m_a}$, $L_0 = 0.10$ m.

% ---------------------------------------------------------------------------
\section{Baseline: the model we augment \quad (Eq.~2)}

The same system with $m_a \to 0$, which recovers the 3-DOF gantry exactly (verified in
the MATLAB). Its matrices are ($s_0$, $J_0$ are $s$, $J$ of Section~1 at $m_a{=}0$):
\begin{equation}
  \small
  M(Y) =
  \begin{bmatrix}
    m_1{+}m_2{+}m_b{+}m_h & s_0 & 0 \\
    s_0 & J_0 & -m_h d \\
    0 & -m_h d & m_h
  \end{bmatrix},
  \quad
  s_0 = (m_1{-}m_2)\tfrac{L_b}{2}{-}m_h Y,
  \quad
  J_0 = J_b{+}J_h{+}(m_1{+}m_2)\tfrac{L_b^2}{4}{+}m_h d^2{+}m_h Y^2 ,
\end{equation}
\begin{equation}
  \small
  C =
  \begin{bmatrix}
    c_{g1}{+}c_{g2} & (c_{g1}{-}c_{g2})\tfrac{L_b}{2} & 0 \\
    (c_{g1}{-}c_{g2})\tfrac{L_b}{2} & c_{b1}{+}c_{b2}{+}(c_{g1}{+}c_{g2})\tfrac{L_b^2}{4} & 0 \\
    0 & 0 & c_y
  \end{bmatrix},
  \quad
  K = \begin{bmatrix} 0 & 0 & 0 \\ 0 & k_{b1}{+}k_{b2} & 0 \\ 0 & 0 & 0 \end{bmatrix}.
\end{equation}
With $M(Y)\ddot q = P^\top u - C\dot q - K q$, implemented as an LPV linear-fractional
representation ($M(Y)^{-1} = N(Y)/d(Y)$) and RK4-discretized, the baseline is
\begin{equation}
  \xb_{k+1} = \fb(\xb_k, u_k)\quad(\text{Eq.\ 2a}),
  \qquad
  \tilde{y}_k = \hb(\xb_k, u_k) = \Cbar\,\xb_k + D_d\,u_k \quad(\text{Eq.\ 2b}).
\end{equation}
Mass is conserved: the baseline uses the full payload $m_h = 10.1$; the truth splits it
as $m_h = m_{h,\mathrm{rigid}} + m_a = 9.09 + 1.01$. Hence (truth $-$ baseline) is only
the absorber.

The output map is frozen first-principles: $C_d = [\,P^\top\ |\ 0_{3\times3}\,]$ (positions
only), $P = \left[\begin{smallmatrix} 1 & 1 & 0\\ \tfrac{L_b}{2} & -\tfrac{L_b}{2} & 0\\ 0 & 0 & 1\end{smallmatrix}\right]$,
$D_d = 0$, normalized to $\Cbar[i,j] = C_d[i,j]\,\sigma_{x,j}/\sigma_{y,i}$.

% ---------------------------------------------------------------------------
\section{Augmentation: dynamic parallel \quad (Eq.~3, Table~1)}

We use the \emph{dynamic parallel} structure of Table~1:
\begin{align}
  \xb_{k+1} &= \fb(\xb_k, u_k) + \fa(\xb_k, \xa_k, u_k) & &\text{(correct the baseline states)}\\
  \xa_{k+1} &= \ga(\xb_k, \xa_k, u_k) & &\text{(drive the added states)}\\
  \yh_k &= \hb(\xb_k, u_k) = \Cbar\,\xb_k + D_d u_k & &(\ha = 0)
\end{align}
The model is a simulation: the encoder sets the initial state
$\xh_{k|k} = \psi(u_{k-n_b:k},\,y_{k-n_a:k})$ (Eq.~5c), the recursion above is rolled out
over the horizon, and $\yh$ is compared to $y$ in the loss (Eq.~5a).

For the gantry, each function and our design choice:
\begin{itemize}
  \item $\fb$: the baseline gantry map (Section~2).
  \item $\fa$: correction to the baseline states, \textbf{nonzero only on the $\Theta$
        components} (rotational position and velocity). Reason: the $X$ and $Y$
        coordinates are pure integrators ($K = 0$), so an additive correction there
        accumulates without a restoring force and diverges.
  \item $\ga$: the dynamics of the two added states $\xa$, targeting $[\delta_a,\dot\delta_a]$.
  \item $\ha = 0$: no learned output term; the output stays $100\%$ baseline.
  \item $\fa,\ga$ are realized by a single feedforward network with $\tanh$ activation.
        $\tanh$ is required because the $Y$-scheduled coupling needs bilinear terms such
        as $Y\!\cdot\!\delta_a$, which a linear (identity) map cannot form.
\end{itemize}
This is a fixed-interconnection special case of the general LFR form (Eq.~4).

% ---------------------------------------------------------------------------
\section{Normalization \quad (Eq.~9)}

Inputs and outputs are scaled to zero mean and unit variance ($T_u$, $T_y$). For the
state scaling $T_x$, Eq.~9 prescribes $\sigma_x$ from the \emph{baseline model's}
simulated states. We instead take
\begin{equation}
  \hl{\sigma_x \ \text{from the data}}:\quad
  x_{\text{log}} = \big[\,P^{-\top} y,\ \ \mathrm{diff}(P^{-\top} y)\,f_s\,\big],
  \qquad \sigma_x = \mathrm{std}(x_{\text{log}}),
\end{equation}
i.e.\ from finite-differenced logical states of the measured (truth) output. The encoder's
linear-model normalization does use baseline-simulated states, so two state-statistics
sources coexist. This is a deviation from Eq.~9 and a candidate baseline-vs-data mismatch
(most likely on the finite-difference velocity scale).

% ---------------------------------------------------------------------------
\section{Training objective and target \quad (Eq.~5)}

Encoder, augmentation, and (optionally) baseline parameters are estimated by minimizing
the truncated multi-step simulation loss
\begin{equation}
  V_{\mathrm{trunc}} = \frac{1}{2N(T{+}1)}\sum_{i=1}^{N}\sum_{\ell=0}^{T-1}
  \big\| \yh_{k_i+\ell\,|\,k_i} - y_{k_i+\ell} \big\|_2^2 .
\end{equation}
Acceptance target: noiseless first, then reach the output-noise floor
$\sigma_n = \mathrm{rms}(y)\,10^{-\mathrm{SNR}/20}$ on the validation simulation RMS.

% ---------------------------------------------------------------------------
\begin{center}
\fbox{\begin{minipage}{0.95\linewidth}
\textbf{Open questions for Jan}
\begin{itemize}
  \item Normalization: $\sigma_x$ from data vs.\ baseline-model states (Eq.~9). Is this a
        problem, and should the physical block and encoder share one source?
  \item Are the added states $\xa$ informative, i.e.\ do they track $\delta_a$ (augment,
        not replace), or does the error drop while $\xa$ stays arbitrary?
  \item Is restricting $\fa$ to the $\Theta$ components adequate, given the true coupling
        also reaches $X$ and $Y$ through $M(Y,\delta_a)$?
\end{itemize}
\end{minipage}}
\end{center}

% ---------------------------------------------------------------------------
\section*{Appendix: parameter values}

\begin{tabular}{@{}llll@{}}
$m_1 = 10.2$ & $m_2 = 10.7$ & $m_b = 22.8$ & $m_h = 10.1$ \\
$J_b = 1.0$ & $J_h = 0.05$ & $L_b = 0.725$ & $d = 0.1$ \\
$c_{g1} = 14.5$ & $c_{g2} = 20.3$ & $c_y = 10$ & $c_{b1}=c_{b2}=9$ \\
$k_{b1}=k_{b2}=1987.5$ & $m_a = 1.01$ & $m_{h,\mathrm{rigid}} = 9.09$ & $L_0 = 0.10$ \\
$f_a = 150$ Hz & $\zeta_a = 0.05$ & $k_a = m_a(2\pi f_a)^2$ & $c_a = 2\zeta_a\sqrt{k_a m_a}$ \\
\end{tabular}
\\[3pt]
{\small Units: masses [kg], inertias [kg\,m$^2$], lengths [m], viscous damping
[N/(m/s)] or [Nm/(rad/s)], stiffness [N/m] or [Nm/rad]. Gantry parameters as in the
Garcia/Jasper baseline.}

\end{document}
